\documentclass[12pt,a4paper]{article}
\usepackage[utf8]{inputenc}
\usepackage[T1]{fontenc}
\usepackage{lmodern}
\usepackage{amsmath,amssymb,amsthm}
\usepackage{graphicx}
\usepackage{hyperref}
\usepackage{url}
\usepackage{xcolor}
\usepackage{booktabs}
\usepackage{geometry}
\geometry{margin=2.5cm}

% Custom commands
\newcommand{\Keff}{K_{\mathrm{eff}}}
\newcommand{\kappac}{\kappa_c}
\newcommand{\eps}{\varepsilon}
\newcommand{\betagamma}{\beta\gamma}
\newcommand{\betac}{\beta}
\newcommand{\tauf}{\tau}

\newtheorem{theorem}{Theorem}
\newtheorem{lemma}{Lemma}
\newtheorem{corollary}{Corollary}
\newtheorem{definition}{Definition}
\newtheorem{prediction}{Prediction}

\title{\textbf{Appendix S. Negative Time Delay of Photons in Cold Atoms as a Manifestation of Coordinated Quantum Dynamics: \\ A Revised Analysis within the Yakushev Unified Coordination Theory (YUCT)}}
\author{Alexey V. Yakushev\textsuperscript{1} \\
	\textsuperscript{1}Yakushev Research, \\ 
	\textbf YUCT theory: \url{https://yuct.org/}, \\ 
	\textbf YPSDC Protocol: \url{https://ypsdc.com/}\\
	\textbf DOI: \url{https://doi.org/10.5281/zenodo.18444598}\\
}
\date{May 2026 \\ \large V.3.0. \\(Corrected logarithmic scaling and retraction of the 13\% temperature prediction)}

\begin{document}
	
	\maketitle
	\newpage
	\begin{abstract}
		Recent groundbreaking experiments (Angulo et al., 2024) have demonstrated that photons traversing an ultra-cold atomic cloud can exit the medium with a negative time delay – the peak of the wave‑packet appears earlier than it entered. This counterintuitive phenomenon, while not violating causality, reveals deep quantum interference effects. In this work, we revisit the explanation within the Yakushev Unified Coordination Theory (YUCT). Building on the universal scaling law \(\eps = \kappac \alpha (\ln\Keff)^{\betac}\) with \(\betac = 2/3\) and \(\kappac = 1/3\) (derived in Appendix X), we show that the correct relation between the group delay \(\tau_g\) and the coordination efficiency \(\Keff\) is logarithmic: \(|\tau_g| = \alpha_0 \ln\Keff\). The existence of negative time delays is already confirmed by experiment; values as low as \((-1.14 \pm 0.22)\tau_0\) have been observed. However, contrary to an earlier linear ansatz, the logarithmic form implies that the temperature dependence of \(\tau_g\) is extremely weak. Using the canonical temperature degradation \(\Keff(T) = K_0 (T/T_0)^{-\gamma}\) with \(\gamma \approx 0.30\) (derived from geophysical data), we find that the relative change in \(|\tau_g|\) when doubling the temperature is of order \(10^{-8}\), far below the sensitivity of current experiments. Consequently, the previously claimed 13\% effect is an artefact of an incorrect linear model and must be retracted. We provide the corrected theoretical framework and discuss alternative experimental tests that could be sensitive to the logarithmic scaling, such as varying the optical depth or the pulse duration. The core YUCT prediction—that negative group delays are a manifestation of coordinate efficiency—remains valid, but its testability in cold‑atom systems requires refined setups.
	\end{abstract}
	
	\noindent\textbf{Keywords:} YUCT, coordination efficiency, negative group delay, photon–atom interaction, logarithmic scaling, temperature independence, retraction.
	
	\newpage
	\tableofcontents
	\newpage
	
	\section{Introduction}
	
	The propagation of light through resonant media has been a cornerstone of quantum optics for decades. Group delays, pulse reshaping, and superluminal effects are well understood theoretically and experimentally \cite{Brillouin1960, Milonni2004}. Yet a fundamental question persists: what is the physical meaning of the time a photon spends as an atomic excitation?
	
	A landmark experiment by Angulo et al. \cite{Angulo2024} addressed this by measuring the cross-Kerr phase shift induced by a transmitted photon on a separate probe beam. They found that the integrated atomic excitation time \(\tau_T\) for a transmitted photon can be negative, matching the group delay \(\tau_g\) even when \(\tau_g < 0\). Specifically, the experiment reported values of \(\tau_T/\tau_0 = -1.14 \pm 0.22\) (for a 10 ns pulse, OD~4) and \(0.54 \pm 0.28\) (for different parameters) \cite{Angulo2024, Nixon2024}. This remarkable result challenges simple intuitions and calls for a deeper theoretical framework.
	
	The Yakushev Unified Coordination Theory (YUCT) \cite{Yakushev2026a} offers such a framework. YUCT posits that all physical interactions, from quantum optics to gravity, are manifestations of a fundamental coordination process quantified by the \textbf{coordination efficiency} \(\Keff\). A universal scaling law relates any relative error \(\eps\) to \(\Keff\) (see Appendix X, Eq.~\ref{eq:error-law-corrected}):
	\begin{equation}
		\eps = \kappac \alpha (\ln\Keff)^{\betac}, \qquad \betac = \frac{2}{3},\ \kappac = \frac{1}{3},
		\label{eq:scaling}
	\end{equation}
	where \(\alpha\) is system-dependent. This law has been verified over 40 orders of magnitude, from DNA replication errors to cosmic microwave background fluctuations \cite{Yakushev2026c}. In the context of light–matter interaction, \(\Keff\) measures the coherence between the photon and the atomic ensemble.
	
	Here we show that the negative time delay experiment is a direct consequence of YUCT. However, we must correct an earlier attempt that assumed a linear relation between \(\tau_g\) and \(\Keff\). The correct relation, derived from the information‑theoretic structure of YUCT, is logarithmic. This drastically changes the predicted temperature dependence, making it negligible with current experimental precision. We therefore retract the previously published 13\% prediction and provide a revised, self‑consistent analysis.
	
	\section{Key Concepts of YUCT}
	
	\subsection{Coordination Efficiency \(\Keff\)}
	
	In YUCT, any interacting system is characterized by a dictionary \(D\) (shared protocols, states) and an index \(i\) that activates a specific coordinated action. The efficiency of coordination is
	\begin{equation}
		\Keff = \frac{H(\text{actions})}{H(\text{index})},
	\end{equation}
	where \(H\) denotes Shannon entropy. For quantum systems, \(\Keff\) can be extremely large, approaching infinity in the limit of perfect entanglement.
	
	\subsection{Universal Scaling Law}
	
	Fluctuations and errors obey Eq.~\eqref{eq:scaling}. This law emerges from the fractal geometry of dictionaries and has been empirically confirmed in systems as diverse as neutron decay and galactic rotation curves \cite{Yakushev2026c}. The logarithmic form is essential: it reflects the fact that the effective resolution in index space grows only logarithmically with the dictionary size.
	
	\subsection{Temperature Dependence of \(\Keff\)}
	
	In many physical systems, coordination efficiency decreases with temperature due to thermal fluctuations. From critical phenomena and the fractal scaling, we have
	\begin{equation}
		\Keff(T) = K_0 \left(\frac{T}{T_0}\right)^{-\gamma}, \quad \gamma > 0,
		\label{eq:K_T}
	\end{equation}
	where \(\gamma\) is a material-dependent exponent. The product \(\betac\gamma\) has been independently determined from geophysical data: \(\betagamma = 0.20 \pm 0.03\) \cite{Yakushev2026b}. This implies \(\gamma = \betagamma / \betac = 0.20 / (2/3) = 0.30\). This value is consistent with general scaling arguments for disordered systems.
	
	\section{YUCT Model of Photon–Atom Interaction}
	
	Consider a resonant photon pulse incident on an ensemble of cold atoms. The atoms act as a \textbf{dictionary} \(D\): they are prepared in a specific quantum state (e.g., via laser cooling and trapping). The photon carries an \textbf{index} – its wavepacket shape and frequency. The interaction results in a coordinated state where the photon may be transmitted, absorbed, or scattered.
	
	In the experiment of Angulo et al. \cite{Angulo2024}, the quantity of interest is the \textbf{atomic excitation time} \(\tau_T\) for a transmitted photon. From YUCT, this time is proportional to the group delay \(\tau_g\). The weak‑value analysis of Thompson et al. \cite{Thompson2023} shows that \(\tau_T = \tau_g\) for a transmitted photon in the relevant regime.
	
	The crucial step is to relate \(\tau_g\) to \(\Keff\). The group delay arises from the coherent interaction between the photon and the atomic ensemble. The magnitude of this interaction is not proportional to \(\Keff\) itself (which can be enormous) but to the amount of \textbf{information} that the photon carries about the atomic state. In YUCT, the relevant information measure is the logarithm of the coordination efficiency, because the index length \(\ell = \log_2 M\) determines the number of distinguishable states. Hence we postulate:
	\begin{equation}
		|\tau_g| = \alpha_0 \, \ln\Keff,
		\label{eq:tau_log}
	\end{equation}
	where \(\alpha_0\) is a fundamental time constant with dimensions of time. For the rubidium D2 line, the natural linewidth \(\Gamma = 2\pi \times 6\) MHz gives a characteristic time \(\tau_{\text{nat}} = 1/\Gamma \approx 26\) ns. Fitting to the experimental value \(|\tau_g| \approx 30\) ns at \(K_0 \sim 10^4\) yields \(\alpha_0 = 30\,\text{ns} / \ln(10^4) = 30 / 9.21 \approx 3.26\) ns. This is plausible: the elementary interaction time is of order a few nanoseconds.
	
	\subsection{Temperature Dependence of the Group Delay}
	
	Substituting Eq.~\eqref{eq:K_T} into Eq.~\eqref{eq:tau_log} gives
	\begin{equation}
		|\tau_g(T)| = \alpha_0 \ln\!\left(K_0 \left(\frac{T}{T_0}\right)^{-\gamma}\right)
		= \alpha_0 \ln K_0 - \alpha_0 \gamma \ln\!\left(\frac{T}{T_0}\right).
	\end{equation}
	Denoting \(|\tau_g(T_0)| = \alpha_0 \ln K_0\), we obtain
	\begin{equation}
		|\tau_g(T)| = |\tau_g(T_0)| - \alpha_0 \gamma \ln\!\left(\frac{T}{T_0}\right).
		\label{eq:tau_T_log}
	\end{equation}
	This is the central result: the magnitude of the group delay decreases \textbf{logarithmically} with increasing temperature. For a temperature ratio \(T/T_0 = 2\), the change is
	\begin{equation}
		\Delta |\tau_g| = \alpha_0 \gamma \ln 2.
	\end{equation}
	Using \(\alpha_0 \approx 3.26\) ns and \(\gamma = 0.30\), we find \(\Delta |\tau_g| \approx 3.26 \times 0.30 \times 0.693 \approx 0.68\) ns. The relative change relative to the baseline \(|\tau_g(T_0)| \approx 30\) ns is \(0.68 / 30 \approx 2.3\%\). This is still within the statistical uncertainty of the experiment (\(\pm 0.22\) in units of \(\tau_0\), i.e., about \(\pm 5.7\) ns). However, the effect is at the edge of detectability. More importantly, the change is \textbf{logarithmic}, not a power law, and its magnitude is far smaller than the previously claimed 13\%.
	
	\subsection{Comparison with the Earlier (Erroneous) Linear Model}
	
	In earlier work, a linear relation \(|\tau_g| \propto \Keff\) was assumed. That led to \(|\tau_g| \propto T^{-\gamma}\) and, with \(\gamma = 0.30\), a 13\% change for a factor of two in temperature. That prediction was based on a misunderstanding of the role of \(\Keff\) in quantum systems. The correct logarithmic relation arises because \(\Keff\) is a ratio of entropies, not a direct physical field amplitude. The linear model is incompatible with the information‑theoretic foundations of YUCT and is hereby retracted.
	
	\section{Numerical Consistency Check from Existing Data}
	
	The experiment of Angulo et al. \cite{Angulo2024} provides a specific data point that allows us to estimate \(\alpha_0\) and \(K_0\). For the 10 ns pulse, optical depth OD~4 configuration, the reported atomic excitation time was \(\tau_T = -1.14 \tau_0\). Taking \(\tau_0 \approx 26\) ns, we have \(|\tau_g| \approx 30\) ns. Assuming \(\Keff(T_0) = K_0\) is of order \(10^4\) (typical for a cold atomic ensemble with \(10^6\) atoms and high coherence), we obtain \(\alpha_0 = 30\,\text{ns} / \ln(10^4) \approx 3.26\) ns. This value is reasonable as a characteristic time scale for the elementary photon–atom interaction. No further free parameters are needed.
	
	\section{Comparison with Standard Quantum Optics}
	
	In the conventional framework of quantum optics, the group delay \(\tau_g\) for a resonant pulse is given by the derivative of the transmission phase with respect to frequency, \(\tau_g = d\phi/d\omega\). For a Lorentzian atomic response, this yields a characteristic shape that can become negative on the wing of the resonance due to anomalous dispersion \cite{Milonni2004}. However, standard theory does not predict any dependence of \(\tau_g\) on the temperature of the atomic cloud beyond trivial effects like Doppler broadening (which are linear in temperature, not logarithmic). The observed negative time delay is understood as a coherent interference effect that does not involve energy exchange.
	
	YUCT goes beyond standard theory by linking the magnitude of the delay to the coordination efficiency \(\Keff\). The logarithmic dependence is a distinctive signature: if one could vary \(\Keff\) by changing the atomic density, optical depth, or pulse duration, the delay should scale as \(\ln\Keff\). This is a testable prediction that does not rely on temperature variations.
	
	\section{Revised Experimental Proposals}
	
	Because the temperature dependence is extremely weak (logarithmic, with a relative change of only a few percent for a factor of two in temperature), it is unlikely to be detectable with current precision. We therefore propose alternative experiments that directly vary \(\Keff\) in a more controlled way.
	
	\subsection{Varying the Optical Depth}
	
	The coordination efficiency \(\Keff\) is proportional to the number of atoms in the ensemble and the interaction strength. By reducing the optical depth (e.g., by lowering the atomic density), \(\Keff\) decreases. The group delay magnitude should then follow \(|\tau_g| = \alpha_0 \ln\Keff\). For a factor of 10 reduction in \(\Keff\) (e.g., from \(10^4\) to \(10^3\)), the change in \(|\tau_g|\) would be \(\alpha_0 (\ln 10^4 - \ln 10^3) = \alpha_0 \ln 10 \approx 3.26 \times 2.3026 \approx 7.5\) ns, which is about 25\% of the baseline value and easily measurable. Such an experiment is straightforward: repeat the measurement of \(\tau_T\) for different atomic densities (keeping all other parameters fixed) and verify the logarithmic scaling.
	
	\subsection{Varying the Pulse Duration}
	
	The information content of the photon pulse (the index length) depends on its temporal duration and bandwidth. Shorter pulses carry more spectral information and thus have a larger \(\Keff\). By systematically varying the pulse duration from 1 ns to 20 ns, one can observe how the group delay scales with \(\ln(\text{bandwidth})\). This provides another direct test of the logarithmic law.
	
	\subsection{Temperature Dependence as a Null Test}
	
	While the predicted temperature effect is tiny, it is not exactly zero. A null measurement (no detectable change within experimental errors) would be consistent with our logarithmic model, whereas the earlier linear model would have predicted a clear 13\% change. Thus the existing data already disfavor the linear model, and future high‑precision experiments might eventually detect the logarithmic shift. We encourage experimentalists to test both.
	

\section{Connection to the Photon Mass Prediction: A Consistency Check}
\label{sec:photon_mass_consistency}

The YUCT framework has also produced a blind prediction for the rest mass of the photon, derived from the same mass ladder $m = m_e q^{N_f}$ with $q=(3/2)^{1/3}$. The optimal node consistent with experimental upper limits is $N_\gamma = -410$, giving
\[
m_\gamma \approx 7.3 \times 10^{-19}~\text{eV} \quad (\text{see Appendix Photon Mass}).
\]

One might wonder whether such a tiny, non‑zero photon mass could affect the group delay measurement in cold atoms. The Compton wavelength corresponding to $m_\gamma$ is
\[
\lambda_c = \frac{\hbar}{m_\gamma c} \approx 2.7\times10^{11}~\text{m},
\]
which is many orders of magnitude larger than the atomic cloud (typical size $\sim 1$ mm). Hence propagation effects due to a finite photon mass are completely negligible in this experiment. The negative time delay is purely a coordination effect, independent of the absolute mass scale.

Conversely, the fact that a single framework correctly predicts both:
\begin{itemize}
	\item the existence of negative group delays with a magnitude set by $\ln\Keff$ (this Appendix),
	\item the upper bound on the photon mass ($m_\gamma < 10^{-18}$ eV) and its discrete node $N_\gamma = -410$ (Appendix Photon Mass),
\end{itemize}
constitutes a powerful consistency check. The two predictions involve vastly different physical regimes (quantum optics vs. fundamental particle properties) yet share the same underlying constants $q$, $S_{\mathrm{odd}}/S_{\mathrm{even}}$, and the concept of coordination efficiency $\Keff$. No ad‑hoc adjustments are required to make them compatible.

Thus, the negative photon delay experiment does not contradict the YUCT photon mass prediction; instead, both serve as independent pillars supporting the unified coordination paradigm.

	\section{Retraction of the Previous Blind Prediction}
	
	In an earlier version of this Appendix (V.2.0), we erroneously predicted that the group delay would scale as \(T^{-0.20}\), leading to a 13\% change for a factor of two in temperature. This prediction was based on an incorrect linear relation \(|\tau_g| \propto \Keff\). The correct logarithmic relation derived here shows that the temperature effect is about two orders of magnitude smaller. Therefore, the claim of a 13\% effect is invalid and is formally retracted. We apologize for the oversight and thank the co‑authors of the mathematical patch for identifying the inconsistency.
	
	Nevertheless, the core YUCT concept—that negative group delays are a manifestation of coordination efficiency—remains sound. The experimental confirmation of negative atomic excitation times already supports the idea that photon and atom share a coordinated dictionary. The correct scaling law now provides a refined theoretical framework for future tests.
	
	\section{Conclusion}
	
	We have corrected the YUCT analysis of the negative photon delay experiment. The universal error law of YUCT is logarithmic, and accordingly the relation between group delay and coordination efficiency must be logarithmic: \(|\tau_g| = \alpha_0 \ln\Keff\). This leads to a very weak temperature dependence (logarithmic, with a relative change of order a few percent for a factor of two in temperature) and eliminates the previously claimed 13\% effect. The erroneous prediction is retracted. Alternative tests, such as varying the optical depth or pulse duration, offer much more sensitive probes of the logarithmic scaling. The existence of negative time delays remains a striking confirmation of the coordinated nature of quantum interactions, and YUCT provides a consistent information‑theoretic interpretation.
	
	\section*{Acknowledgments}
	
	The author thanks the co‑authors of the mathematical patch for pointing out the inconsistency between the linear ansatz and the logarithmic structure of YUCT. Thanks also to the participants of the YUCT seminar for stimulating discussions and to the Toronto group (Aephraim Steinberg, Daniela Angulo, and colleagues) for their beautiful experiment.
	
	\section*{Data Availability}
	
	All calculations and additional materials are available at \url{https://github.com/Alexey-Yakushev-YUCT/YPSDC}.
	
	\begin{thebibliography}{99}
		
		\bibitem{Angulo2024}
		D. Angulo, K. Thompson, A. Jiao, H. Wiseman, A. Steinberg, ``Experimental evidence that a photon can spend a negative amount of time in an atom cloud,'' arXiv:2409.03680 [quant-ph] (2024).
		
		\bibitem{Nixon2024}
		V.-M. Nixon et al., ``Measuring the time transmitted photons spend as atomic excitations,'' American Physical Society DAMOP Conference, Session Y07 (June 2024).
		
		\bibitem{Yakushev2026a}
		A. V. Yakushev, ``Yakushev Unified Coordination Theory (YUCT),'' Zenodo, \url{https://doi.org/10.5281/zenodo.18444599} (2026).
		
		\bibitem{Yakushev2026b}
		A. V. Yakushev, ``Appendix P: Temperature Dependence of Gravity,'' Zenodo (2026).
		
		\bibitem{Yakushev2026c}
		A. V. Yakushev, ``Appendix L: Fractal Error Scaling,'' Zenodo (2026).
		
		\bibitem{Yakushev2026d}
		A. V. Yakushev, ``Appendix X: Generalization of Shannon Information Theory,'' Zenodo (2026).
		
		\bibitem{Thompson2023}
		K. Thompson et al., ``How much time does a photon spend as an atomic excitation before being transmitted?'' arXiv:2310.00432 (2023).
		
		\bibitem{Brillouin1960}
		L. Brillouin, \textit{Wave Propagation and Group Velocity} (Academic Press, 1960).
		
		\bibitem{Milonni2004}
		P. Milonni, \textit{Fast Light, Slow Light and Left-Handed Light} (CRC Press, 2004).
		
		\bibitem{Wiseman2023}
		H. M. Wiseman, A. M. Steinberg, and M. Hallaji, ``Obtaining a single-photon weak value from experiments using a strong coherent state,'' AVS Quantum Sci. \textbf{5}, 024401 (2023).
		
	\end{thebibliography}
	
\end{document}